
% ============================================================================
\clearpage
\section*{v4.10.0-r27 Source-Only: Correction to r26, the Image Audit of HS(6) Is in the Atlas}
\addcontentsline{toc}{section}{v4.10.0-r27 Source-Only: Correction to r26, the Image Audit of HS(6) Is in the Atlas}
\label{sec:v4100-r27-image-audit-correction}
% ============================================================================

\paragraph{Purpose.}
v4.10.0-r26 recorded the image audit of the HS(6) coding map, which the atlas layer r58 recommends before any new shift, as open.  The
author asked for that audit to be carried out (2026-11-16).  Reading the atlas before computing anything shows that the audit is already
there, in the layer directly after r58.  This layer corrects r26, records what the atlas states, and checks the atlas's reasons against later work.  No computation, no new result.  The
notes are \texttt{notes/\allowbreak research/\allowbreak v4\_10\_0\_r27/}.  Earlier layers are not rewritten.

\noindent\textbf{Source-only status.}
This v4.10.0-r27 layer is source-only.  It does not build a PDF, does not assign a new DOI,
and does not authorize a public release.  The v5.1.0 Version DOI
\href{https://doi.org/10.5281/zenodo.24237112}{10.5281/zenodo.24237112} and the Concept DOI
\href{https://doi.org/10.5281/zenodo.19210728}{10.5281/zenodo.19210728} are unchanged by it.

\subsection*{v4.10.0-r27 results}
\label{subsec:v4100-r27-results}

\paragraph{Correction to v4.10.0-r26.3.}
r26.3 and the register of r26 call the image audit of HS(6) open.  That is wrong.  The atlas layer r59 (v470/32, ``Dixon image audit'')
carries it out: it classifies the reachable image of the coding map over the capture classes \(A_1,\dots,A_6\) as Dixon-admissible (\(A_1\);
\(A_6\) routed or admissible), routed (\(A_2\), \(A_3\)), blocked (\(A_4\)) and HC/MM absorbed (\(A_5\)), and proves that no residual-open class
remains, under the zero-register hypotheses (HS(4) leakage, HS(5) debt, hidden leakage and HC/MM debt equal to zero) and exhaustiveness of
the capture map.  The recommendation of r58 is thereby fulfilled.  The other statements of r26 (Definition r26.1, Theoremlet r26.2, the
second limit of r26.3 on the change of meaning) are unaffected.

\paragraph{Remark v4.10.0-r27.1 (what the atlas leaves open after r59).}
The layers after r59 audit the route menu.  The integrated residual register concludes that the local route image is exhausted and that
OP5 is locally closed inside the declared delta-stack grammar; it remains open only at the full Active-Core integration and exhaustiveness
boundary, \(\mathcal R^{\mathrm{loc}}_{\mathrm{mono}}=\varnothing\) still to be shown.  Two further layers (monolithic promotion audit, closure
hardening and old-branch mapping) continue this.  The atlas also states, in r59, that the next useful layer ``should not add a new algebraic
carrier''.  HS(7) and HS(8) of r25 and r26 are such carriers; their use as generation carriers lies outside the purpose of the atlas
tower, as r26.3 already said.

\paragraph{Remark v4.10.0-r27.2 (why the atlas advised against new shifts, and what of it still holds).}
On the author's request (2026-11-16) the reasons behind the advice of r58 (``the next meaningful layer is not a new shift'') and r59
(``should not add a new algebraic carrier'') were checked against what happened since.  The atlas gives three.
\begin{enumerate}[leftmargin=2em]
\item[(a)] \emph{A finite OP5 audit.}  The Dixon coding turned the OP5 tunnel audit ``into a finite Dixon-mask admissibility audit rather
than an unbounded route search''; a new carrier would reopen the route grammar.  Since then the old OP5 question was suspended and rebuilt:
the historical \(\Lambda\)-cancellation duty is no established consequence of the core, and the source--return package is a conditional
schema without a native route (v4.10.0-r05, r11.2(a)).  The local OP5 closure itself contains no frame quantity and stands within its
grammar (v4.9.0-r23(f)).  Reason (a) still holds for coding OP5 routes; it does not bind other uses of the carrier.
\item[(b)] \emph{Algebraic caution.}  HS(4) is ``the first level at which a naive 16-channel extension may behave like a sedenion-type
expansion''; the atlas therefore gates it by the projector and a leakage register instead of using sedenions as a free algebra.  The
concern is the Cayley--Dickson doubling of the algebra, beyond the octonions.  The rule of r26.1 doubles the carrier into copies, as HS(5)
does, and doubles no algebra; the copies stay \(\C\otimes\Hh\otimes\Oo\).  Reason (b) does not apply to HS(7) and HS(8) as copies.  It would
apply to a sedenion reading of 128 channels, route~A of r13.
\item[(c)] \emph{Scope.}  The Dixon carrier is used as a route code, with no Standard-Model claim.  Since v4.8 the project uses
\(\C\otimes\Hh\otimes\Oo\) explicitly as the carrier of one generation (CT-003, r13, r16.1, r19, r25).  Both roles now coexist; reason (c) is a
scope statement for the atlas layer, not a theorem against the second role.  It asks that the two roles be kept apart: route masks are not
generation states.
\end{enumerate}
Reading: the advice was a scope decision for the OP5 programme of v4.7, plus a caution against algebra doubling.  Neither forbids HS(7)
and HS(8) as copy carriers for generations.  For OP5 route coding the advice stays sensible.  The atlas works in its ``photon frame'',
typed R by v4.9.0-r24; the copy index of r25/r26 is frame-free, so no frame conflict arises.

\paragraph{What this does not give.}
No new audit and no computation; the content is that of the atlas.  The image audit of r59 is about OP5 routes on HS(6); it says nothing
about HS(7), HS(8) or the generation reading of r25.

\subsection*{v4.10.0-r27 status in the three frames}
\label{subsec:v4100-r27-frames}

\begin{itemize}[leftmargin=2em]
\item M, T, P: no frame statement; the atlas audit is set in the photon-frame HC-halo grammar of the local proof chain.
\end{itemize}

\subsection*{v4.10.0-r27 register entries}
\label{subsec:v4100-r27-register}

\begin{center}
\small
\begin{tabular}{@{}>{\raggedright\arraybackslash}p{0.34\textwidth}>{\raggedright\arraybackslash}p{0.58\textwidth}@{}}
\toprule
statement & status after v4.10.0-r27\\
\midrule
image audit of HS(6) & done in atlas layer r59; no residual-open class under the zero-register and exhaustiveness hypotheses; r26.3 corrected\\
OP5 (atlas) & locally closed inside the declared grammar; open at the full Active-Core exhaustiveness boundary\\
HS(7), HS(8) & admissible under the setting r26.1; outside the purpose of the atlas tower\\
reasons of the atlas against new shifts & OP5 audit finiteness (still valid for routes), caution against algebra doubling (does not apply to copies), scope (two roles to be kept apart) (r27.2)\\
\bottomrule
\end{tabular}
\end{center}

\subsection*{v4.10.0-r27 ledger and guard update}
\label{subsec:v4100-r27-ledger-guard}

\[
R_{\rm chk}^{v4100\text{-}r27}=(1,0,0)
\]
for (correction recorded; no computation; nothing new opened).

The layer spends no guard:
\[
R_{\rm glob}^{(v4100\text{-}r27,+)}=(0,1,1,1,1),
\qquad
\Gamma_{\rm guard}^{v4100\text{-}r27}=(1,1,1,1,1),
\qquad
\beta_4=1/256\ \text{unspent}.
\]
The layer does not imply exact finite \(C5\) zero closure, exact all-readout terminality, \(C6/T2\), a global photon-clock/RH theorem,
or global all-readout closure.
